\documentclass[11pt,a4paper,fontset=fandol]{ctexart}

\usepackage[margin=1in]{geometry}
\usepackage{amsmath,amssymb,bm}
\usepackage{booktabs}
\usepackage{longtable}
\usepackage{array}
\usepackage{enumitem}
\usepackage{xcolor}
\usepackage{hyperref}
\usepackage{listings}
\usepackage{fancyhdr}
\usepackage{titlesec}

\hypersetup{
  colorlinks=true,
  linkcolor=blue,
  urlcolor=blue,
  citecolor=blue
}

\definecolor{codegray}{RGB}{245,245,245}
\definecolor{darkblue}{RGB}{20,60,120}
\definecolor{darkred}{RGB}{150,40,40}
\definecolor{darkgreen}{RGB}{30,120,60}

\lstset{
  backgroundcolor=\color{codegray},
  basicstyle=\ttfamily\small,
  breaklines=true,
  frame=single,
  columns=fullflexible,
  keepspaces=true,
  showstringspaces=false
}

\pagestyle{fancy}
\fancyhf{}
\lhead{AI-assisted LQCD Meson DA Agent}
\rhead{Month 1 Guide}
\cfoot{\thepage}

\titleformat{\section}{\Large\bfseries\color{darkblue}}{\thesection}{0.8em}{}
\titleformat{\subsection}{\large\bfseries\color{darkblue}}{\thesubsection}{0.8em}{}
\titleformat{\subsubsection}{\normalsize\bfseries\color{darkblue}}{\thesubsubsection}{0.8em}{}

\title{\bfseries 第一个月详细指导书\\AI-assisted LQCD Meson DA Analysis Agent}
\author{项目训练材料：面向研究生的学习、分析与搭建路线}
\date{版本：Month 1 Detailed Guide}

\begin{document}
\maketitle
\tableofcontents
\newpage

\section{第一个月的定位}

第一个月的核心目标不是立刻做出完整的高级 agent 系统，而是建立一个\textbf{可靠、可运行、可测试、可复现、可审计}的基础版本。学生需要通过第一个月完成从“会运行别人代码”到“能把科研分析流程工程化”的转变。

第一个月结束时，学生应该能够做到：

\begin{enumerate}[leftmargin=2em]
  \item 复现一个已经准备好的 golden example，例如 \lstinline|examples/l64c64a076_m140| 或导师指定的 meson DA 分析样例；
  \item 说清楚这个样例的输入、输出、中间文件、关键图和主要物理检查；
  \item 把原本 notebook-style 或手动运行的流程整理成 config-driven Python CLI；
  \item 使用 Pydantic 定义 config schema 和 result schema；
  \item 使用 pytest 对关键模块做基本测试；
  \item 使用 Codex 辅助写代码、修 bug、补 tests，但不盲目信任 Codex；
  \item 用 Git 做小步提交，保证每一步可以回退；
  \item 输出一个基础的 run manifest、result.json、summary.csv、plots 和 run.log；
  \item 写出第一版项目说明和复现实验报告。
\end{enumerate}

第一个月不要求学生完全掌握 Snakemake 和 LangGraph，但需要理解它们后续的定位：

\begin{itemize}[leftmargin=2em]
  \item Snakemake 后续用于管理真实分析流程的文件依赖和增量重跑；
  \item LangGraph 后续用于组织多 agent 状态流转、human approval、targeted rerun 和 memory update；
  \item 第一个月的重点是先把每个基本模块、输入输出和测试做清楚，否则后续 workflow engine 和 agent orchestration 会变成空架子。
\end{itemize}

\section{第一个月的总体交付物}

第一个月结束时，学生应提交以下内容：

\begin{lstlisting}
lqcd-meson-da-agent/
  README.md
  docs/
    month1_reproduction_report.md
    data_flow_map.md
    codex_usage_log.md
  configs/
    golden_example.yaml
  src/
    lqcd_agent/
      __init__.py
      schema/
        analysis_config.py
        analysis_result.py
      audit/
        data_audit.py
      cli/
        run_golden.py
      utils/
        logging_utils.py
        io_utils.py
  tests/
    test_schema.py
    test_data_audit.py
    test_cli_smoke.py
  outputs/
    golden_example/
      manifest.json
      result.json
      summary.csv
      plots/
      run.log
\end{lstlisting}

如果当前项目已有不同目录结构，学生不需要强行照搬上述结构，但必须满足以下原则：

\begin{enumerate}[leftmargin=2em]
  \item 输入配置集中管理；
  \item 原始数据不被修改；
  \item 输出目录有明确 run id；
  \item 每次运行留下 manifest 和 log；
  \item 关键结果以 JSON/CSV 结构化保存；
  \item 至少有最小 pytest 测试；
  \item README 能让另一个人照着复现。
\end{enumerate}

\section{学生每天的工作方式}

学生每天都应按照下面的循环工作。

\subsection{Step 1：先写当天小目标}

不要一上来让 Codex “帮我做完整项目”。每天先写一个小目标，例如：

\begin{lstlisting}
Today goal:
Run the golden example from a clean environment, record all required inputs and outputs, and write a data-flow map.
\end{lstlisting}

或者：

\begin{lstlisting}
Today goal:
Add a Pydantic schema for the golden example config and write tests for missing fields and invalid paths.
\end{lstlisting}

\subsection{Step 2：让 Codex 做小而具体的任务}

Codex prompt 必须包含：

\begin{itemize}[leftmargin=2em]
  \item 当前项目背景；
  \item 要修改哪些文件；
  \item 输入和输出是什么；
  \item 不允许修改什么；
  \item 需要添加哪些 tests；
  \item 完成后如何运行检查。
\end{itemize}

\subsection{Step 3：学生自己运行和检查}

Codex 写完后，学生必须自己运行：

\begin{lstlisting}[language=bash]
pytest
python -m lqcd_agent.cli.run_golden --config configs/golden_example.yaml
\end{lstlisting}

同时检查输出文件是否真的生成，图是否能打开，JSON/CSV 是否符合预期。

\subsection{Step 4：用 traceback 让 Codex 修复}

如果报错，不要简单说“修一下”。应该把 traceback、运行命令、预期行为和实际行为给 Codex。

\begin{lstlisting}
The command failed:
python -m lqcd_agent.cli.run_golden --config configs/golden_example.yaml

Traceback:
[paste traceback]

Expected behavior:
The CLI should load the config, run the golden example, and write manifest.json.

Please fix only the cause of this error. Do not refactor unrelated files.
Also add a regression test if appropriate.
\end{lstlisting}

\subsection{Step 5：Git 小步提交}

每个完成的小任务都应该提交一次。

\begin{lstlisting}[language=bash]
git status
git diff
git add <changed files>
git commit -m "Add config schema for golden example"
\end{lstlisting}

不要把一周所有修改堆到一个 commit。

\section{Week 1：真实 example 与统计物理基础}

这一节是 Month 1 的核心背景。它把真实 example 的数据流、统计直觉和代码结构串在一起，目标不是先背概念，而是先把“看图、动手、复现”这一条最小闭环跑通。

Week 1 的目标不是让学生“看懂结果”，而是让学生建立最基本的科研判断方式。对于两点函数分析来说，最重要的第一步不是拟合，而是学会读图：知道什么叫 plateau，知道什么时候噪声开始压过信号，知道为什么不同动量的 correlator 不能直接用同一套直觉去解释。

从教学上说，Week 1 的前半段就是一个最小闭环：

\begin{enumerate}[leftmargin=2em]
  \item 读入一组两点 correlator；
  \item 只计算 effective mass；
  \item 只画图，不做复杂拟合；
  \item 从图上回答“这里是不是已经接近单态区”；
  \item 把观察到的 plateau 和噪声结构带到后面的 \lstinline|1-c2pt-fit|。
\end{enumerate}

如果这一步没有学透，后面的 two-point fit 和 DA ratio 会变成纯粹的黑箱操作。

\subsection{Week 1 的学习-练习循环}

\subsubsection{建议阅读路线}

如果你是第一次接触这套分析，建议按下面的顺序读本章，而不是从头到尾随便翻。这里的原则是：先建立统计和物理直觉，再做最小练习，最后才进入真正的拟合。

\begin{enumerate}[leftmargin=2em]
  \item 先看 \textbf{当前 example 的真实主线}，理解数据从哪里来、要去哪里；
  \item 接着看 \textbf{两点 correlator 的物理模型}，弄清楚 $A_n$ 和 $E_n$ 是什么，以及为什么拟合参数要满足物理约束；
  \item 然后看 \textbf{bootstrap、binning、协方差、$\chi^2$}，理解误差条、相关性和拟合优度是怎么从样本里来的；
  \item 再打开 \textbf{0-effective-mass}，把前面学到的概念立刻用在图上，先建立 plateau、噪声和 bootstrap 的直觉；
  \item 再看 \textbf{1-c2pt-fit}，把 effective mass 的直觉转成真正的多态拟合窗口选择；
  \item 之后看 \textbf{DA ratio 的最小物理图像}，把两点拟合和后面的矩阵元提取连起来；
  \item 最后看 \textbf{notebook 与 input file 的运行方式}，把理解过的流程固化成可复现输入。
\end{enumerate}

这个顺序的原则很简单：先看“数据流”，再看“模型”，再看“统计”，然后立刻用 \lstinline|0-effective-mass| 把概念落到图上，最后再把图上的直觉带进真正的拟合。

更适合初学者的方式，是把每一个知识点都立刻对应到一个动作上：

\begin{enumerate}[leftmargin=2em]
  \item 先看当前 example 的主线和两点 correlator 的模型；
  \item 再理解 bootstrap、binning、协方差和 $\chi^2$ 的统计含义；
  \item 立刻跑 \lstinline|0-effective-mass| 的 notebook，把图画出来；
  \item 立刻再跑 input file，把 notebook 里的理解冻结成可复现输入；
  \item 立刻写下自己观察到的 plateau、噪声和误差变化；
  \item 再把这些观察带到 \lstinline|1-c2pt-fit| 里去选窗口。
\end{enumerate}

\subsection{Week 1 与后续分析的衔接}

\subsubsection{当前 example 的真实主线}

本仓库当前的实际 example 是 \lstinline|examples/l64c64a076_m140/|。它是一条已经裁剪成 DA-only 目标的真实分析链，包含 effective mass、两点拟合、DA bare matrix element、normalization 和 Fourier 变换。

最常见的分析入口是：

\begin{lstlisting}[language=bash]
python -m lqcd_analysis.cli 2pt-effective-mass ...
python -m lqcd_analysis.cli 2pt-nstate-fit ...
python -m lqcd_analysis.cli da-nstate-fit ...
python -m lqcd_analysis.cli da-normalize ...
python -m lqcd_analysis.cli da-fourier ...
\end{lstlisting}

学生至少要理解这四段之间的依赖关系：

\begin{center}
\begin{tabular}{p{0.18\linewidth}p{0.32\linewidth}p{0.40\linewidth}}
\toprule
阶段 & 结果层级 & 关键物理问题 \\
\midrule
0-effective-mass & effective mass 和图 & plateau、噪声、bootstrap 误差是否合理 \\
1-c2pt-fit & 两点 correlator 拟合 & 能量、overlap、$\chi^2/\mathrm{dof}$ 是否稳定 \\
2-bm & DA bare matrix element & 是否已把 ratio 中的矩阵元提取出来 \\
3-normalize & 归一化后的中间量 & 只保留 mode3 路径，检查归一化是否合理 \\
4-FT & Fourier 变换结果 & 最终分布的形状和物理趋势是否正常 \\
\bottomrule
\end{tabular}
\end{center}

当前 example 的数据已经被裁剪成一条更窄但更清楚的 DA-only 主线：输入里只保留 \lstinline|b_X|、\lstinline|OT5| 和 \lstinline|bT = 0|，同时删掉了 \lstinline|b_Y| 和 \lstinline|OZ5|，而下游 analysis 也只保留 \lstinline|mode3| 路径。学生在读这一章时，只需要知道这套 example 的数据长什么样、为什么只剩这些分支，以及这些分支在后续分析里分别扮演什么角色，不需要再记一串 bullet。

这份 Month 1 章节已经把 two-point fit 和 DA ratio-fit 的关键物理、统计和代码逻辑合并在一起了，所以学生不需要再翻独立的旧 guide。真正需要反复看的，是当前 example 的目录结构和输入文件约定。

\subsubsection{两点 correlator 的物理模型}

对任意插入算符 $O$，两点函数的定义是

\begin{equation}
C(t) = \langle O(t)\, O^\dagger(0) \rangle.
\end{equation}

在有限时长 $N_t$ 的格点上，当前实现支持三种模型：

\begin{align}
C_{\mathrm{normal}}(t) &= \sum_{n=0}^{N_{\mathrm{state}}-1} A_n e^{-E_n t}, \\
C_{\mathrm{symmetric}}(t) &= \sum_{n=0}^{N_{\mathrm{state}}-1} A_n \left(e^{-E_n t} + e^{-E_n (N_t - t)}\right), \\
C_{\mathrm{antisymmetric}}(t) &= \sum_{n=0}^{N_{\mathrm{state}}-1} A_n \left(e^{-E_n t} - e^{-E_n (N_t - t)}\right).
\end{align}

这里的 $A_n$ 是 overlap/amplitude，$E_n$ 是能量。代码中不会直接拟合任意无序参数，而是通过对数与 gap 参数化强制：

\begin{itemize}[leftmargin=2em]
  \item $A_n > 0$；
  \item $E_0 > 0$；
  \item $E_{n+1} > E_n$。
\end{itemize}

所以拟合不是“猜一个数”，而是在一个满足物理约束的参数空间里做非线性最小二乘。

\subsubsection{effective mass 的含义}

对于 \lstinline|model = normal|，effective mass 定义为

\begin{equation}
m_{\mathrm{eff}}(t) = \log\frac{C(t)}{C(t+1)}.
\end{equation}

对于 periodic / antiperiodic 的情况，代码会数值求解 cosh / sinh 版本的方程：

\begin{align}
\frac{C(t)}{C(t+1)} &=
\frac{\cosh\!\left[m_{\mathrm{eff}}(t+1 - N_t/2)\right]}
\cosh\!\left[m_{\mathrm{eff}}(t - N_t/2)\right] \quad \text{(symmetric)}, \\
\frac{C(t)}{C(t+1)} &=
\frac{\sinh\!\left[m_{\mathrm{eff}}(t+1 - N_t/2)\right]}
\sinh\!\left[m_{\mathrm{eff}}(t - N_t/2)\right] \quad \text{(antisymmetric)}.
\end{align}

effective mass 的作用不是最终物理结论，而是 plateau 诊断：它帮助判断拟合窗口是否已经进入近似单态区。

\subsubsection{bootstrap、binning 和协方差}

先回答一个最根本的问题：\textbf{为什么这里要 bootstrap sampling？}

因为格点数据里的“误差”不是凭空给出的，而是要从一批带噪声的 configuration 里估计出来。我们关心的不是某一次测量值，而是“如果重新抽一批样本，均值、误差条和拟合参数会怎么波动”。bootstrap 正是在做这件事。

当前代码的统计顺序是：

\begin{enumerate}[leftmargin=2em]
  \item 先对原始 configuration 沿第一维做 binning，把相邻配置压成 block；
  \item 再从 binned blocks 中有放回抽样；
  \item 每个 bootstrap sample 都是一条新的平均 correlator；
  \item 用 bootstrap sample 的分布估计误差、协方差和参数不确定性。
\end{enumerate}

这一步背后的统计含义是：我们并不把单个 configuration 当成“真值”，而是把一批经过 binning 之后的 block 当成经验样本，再用 bootstrap 去逼近“重复做一次测量时，结果会怎样波动”。对于格点数据来说，这一点尤其重要，因为相邻 configuration 可能存在自相关。如果不先 binning，就可能把相关样本误当成独立样本，从而把误差条估得过于乐观。

把这件事翻译成可操作的算法，就是：

\begin{enumerate}[leftmargin=2em]
  \item 从原始 correlator 读入一个形状类似 \lstinline|(n_cfg, Nt)| 的数组；
  \item 按 \lstinline|binsize| 在 configuration 维度上分块平均；
  \item 每次 bootstrap resample 都从这些 blocks 里“有放回”抽取 \lstinline|n_draws| 个 block；
  \item 对抽到的 block 再求平均，得到一条 bootstrap mean correlator；
  \item 重复很多次，最后得到 bootstrap 分布。
\end{enumerate}

对学生来说，bootstrap、binning 和协方差分别回答的是三件不同的事：

\begin{itemize}[leftmargin=2em]
  \item \textbf{binning} 负责降低自相关的影响；
  \item \textbf{bootstrap} 负责把样本波动传播到最终物理量；
  \item \textbf{covariance} 负责告诉你不同时间点是不是一起涨、一起跌。
\end{itemize}

所以这一步不是“为了多算几遍”，而是为了把统计不确定性和时间相关结构都显式地保留下来。

如果再往代码里看一层，真实顺序就是 \lstinline|apply_fold_t| 先利用时间反向对称性把 correlator 整理得更平滑，然后 \lstinline|bin_correlators| 压低自相关，接着 \lstinline|bootstrap_correlator_means| 生成重采样平均值，最后这些 bootstrap means 被送进 \lstinline|effective_mass_with_bootstrap|、协方差估计和拟合模块。

如果第 $b$ 个 bootstrap 样本的 correlator 是 $C^{(b)}(t)$，那么均值可写为

\begin{equation}
\bar{C}(t) = \frac{1}{N_{\mathrm{boot}}}\sum_{b=1}^{N_{\mathrm{boot}}} C^{(b)}(t).
\end{equation}

协方差矩阵由 bootstrap 样本估计：

\begin{equation}
\mathrm{Cov}_{ij} = \frac{1}{N_{\mathrm{boot}}-1}
\sum_{b=1}^{N_{\mathrm{boot}}}
\left(C_i^{(b)} - \bar{C}_i\right)\left(C_j^{(b)} - \bar{C}_j\right).
\end{equation}

如果协方差矩阵不稳定，代码会做 shrinkage：

\begin{equation}
C_\lambda = (1-\lambda)C + \lambda\,\mathrm{diag}(C).
\end{equation}

其中 $\lambda=0$ 保留完整协方差，$\lambda=1$ 只保留对角项。

在本仓库里，bootstrap 的实现顺序是固定的。先把 correlator 组织成 \lstinline|(n_cfg, Nt)| 的数组；然后按 \lstinline|binsize| 沿第一个维度做 binning，得到更稳定的 block；接着用 \lstinline|bootstrap_indices| 在这些 block 上做有放回抽样；最后用 \lstinline|bootstrap_means| 计算每个 resample 的平均 correlator。

这条链可以用下面的简化写法理解：

\begin{lstlisting}[language=Python]
binned = bin_correlators(correlators, binsize=binsize)
indices = bootstrap_indices(n_blocks, draw_size, seed=seed, n_boot=n_boot)
resampled_means = bootstrap_means(binned, indices=indices)
\end{lstlisting}

这里的关键不是“把同一份数据重复算很多次”，而是“把可能相关的 configuration 先压成 block，再把 block 当成经验样本重新抽样”。这样得到的 bootstrap 分布，才更适合用来估计误差、协方差和后续拟合参数的不确定性。对于格点数据来说，这一步尤其重要，因为相邻 configuration 往往并不独立。

如果把这段逻辑写成更接近真实代码的顺序，就是 \lstinline|apply_fold_t| 先利用时间反向对称性把 correlator 整理得更平滑，然后 \lstinline|bin_correlators| 压低自相关，接着 \lstinline|bootstrap_correlator_means| 生成重采样平均值，最后这些 bootstrap means 被送进 \lstinline|effective_mass_with_bootstrap|、协方差估计和拟合模块。

所以 bootstrap、binning 和协方差分别负责三件不同的事：

\begin{itemize}[leftmargin=2em]
  \item binning 减弱自相关；
  \item bootstrap 把样本波动传播到物理量；
  \item covariance 告诉我们不同时间点之间是不是一起涨、一起跌。
\end{itemize}

\subsubsection{最小 \texorpdfstring{$\chi^2$}{chi2} 非线性拟合}

当前 two-point 和 DA 拟合都采用最小 $\chi^2$ 的非线性最小二乘。对数据向量 $d$ 和模型向量 $m(\theta)$，残差写成

\begin{equation}
r(\theta) = \frac{m(\theta)-d}{\sigma},
\end{equation}

或者在 correlated fit 时写成

\begin{equation}
r(\theta) = L^{-1}\!\left(m(\theta)-d\right),
\end{equation}

其中 $L$ 是协方差矩阵的 Cholesky 因子，满足 $C = LL^\top$。

目标函数是

\begin{equation}
\chi^2(\theta) = r(\theta)^\top r(\theta),
\qquad
\chi^2/\mathrm{dof} = \frac{\chi^2(\theta)}{N_{\mathrm{data}} - N_{\mathrm{param}}}.
\end{equation}

如果把统计假设写完整一点，就是：在近似高斯误差下，最大化似然函数等价于最小化 $\chi^2$。所以这里的 fit 不是“凭感觉找一条最像的曲线”，而是一个有明确统计意义的参数估计问题。

代码使用 \lstinline|scipy.optimize.least_squares(method="trf")| 来最小化这个目标。对学生来说，最重要的是理解：

\begin{itemize}[leftmargin=2em]
  \item 小的 $\chi^2/\mathrm{dof}$ 说明模型和窗口与数据较一致，但不自动等于“物理正确”；
  \item 过大的 $\chi^2/\mathrm{dof}$ 往往表示窗口、模型或协方差处理有问题；
  \item correlated fit 更接近真实统计结构，但更依赖稳定协方差；
  \item bootstrap fit 让结果依赖整条重采样分布，而不是某一个单次最优点。
\end{itemize}

在代码里，\lstinline|least_squares| 实际上最小化的是一串标准化残差。先用当前参数构造模型值 \lstinline|m(theta)|，再与数据向量 \lstinline|d| 相减得到 \lstinline|delta = m(theta) - d|。如果是 uncorrelated fit，就把 \lstinline|delta| 除以逐点误差 \lstinline|sigma|；如果是 correlated fit，就先对 bootstrap 协方差做 Cholesky 分解，用 \lstinline|L^{-1} delta| 作为残差。然后把这些残差平方求和，就是 \lstinline|chi^2|。

为了让优化稳定，代码还会给参数做物理约束重参数化：

\begin{itemize}[leftmargin=2em]
  \item amplitude 用 log 参数化，保证它们始终为正；
  \item energy 按 state 顺序写成正的基态能量和正的 gap，从而保证能级递增。
\end{itemize}

换成公式就是：

\begin{equation}
r_i(\theta) = \frac{m_i(\theta)-d_i}{\sigma_i}
\quad \text{或} \quad
r(\theta) = L^{-1}(m(\theta)-d),
\end{equation}

\begin{equation}
\chi^2(\theta)=\sum_i r_i(\theta)^2
\quad \text{或} \quad
\chi^2(\theta)=r(\theta)^\top r(\theta).
\end{equation}

这也是为什么它叫 chi-square fit：它不是在直接最小化“模型和数据的差”，而是在最小化“按照统计不确定度标准化后的差”。只有当偏差相对于误差足够小，这个偏差才会被认为可以接受。

对学生来说，读 chi-square fit 结果时至少要同时看四件事：

\begin{enumerate}[leftmargin=2em]
  \item 参数值是不是落在物理上合理的区间；
  \item \lstinline|chi^2/\mathrm{dof}| 是否说明窗口和模型能接受这批数据；
  \item 协方差是否稳定，是否需要 shrinkage；
  \item bootstrap 后的参数分布是否和 mean fit 一致。
\end{enumerate}

更具体地说，学生读拟合结果时应该同时问四个问题：

\begin{enumerate}[leftmargin=2em]
  \item 这个模型是否真的符合该时间窗里的物理假设？
  \item 当前窗口是否已经尽量避开 excited-state contamination？
  \item 协方差矩阵是否稳定到足以支持 correlated fit？
  \item 当前的 $\chi^2/\mathrm{dof}$ 是“模型足够好”，还是“窗口太短、自由度太少”造成的表面好看？
\end{enumerate}

也就是说，$\chi^2$ 不是一个孤立的数字，而是和窗口选择、模型选择以及协方差处理绑在一起看的。

\subsubsection{bootstrap-centered 的中心值和误差}

当前实现并不是只做一次拟合就结束，而是采用 bootstrap-centered 的汇总方式：

\begin{enumerate}[leftmargin=2em]
  \item 先对 bootstrap mean 做 mean fit，得到 warm start；
  \item 再对每个 bootstrap sample 做拟合；
  \item 对所有成功样本的参数进行稳健汇总；
  \item 用分位数而不是单次 Hessian 误差定义最终中心值和误差。
\end{enumerate}

具体来说，当前代码用 16\% 和 84\% 分位数定义

\begin{equation}
\mu = \frac{P_{16}+P_{84}}{2},
\qquad
\sigma = \frac{P_{84}-P_{16}}{2}.
\end{equation}

所以很多表格里的 central value / uncertainty，都是 bootstrap 分布的分位数摘要。

\subsubsection{从 \texorpdfstring{0-effective-mass}{0-effective-mass} 到 \texorpdfstring{1-c2pt-fit}{1-c2pt-fit}}

这一段是 Week 1 里最重要的过渡：学生不是在学两个孤立步骤，而是在学同一条物理链条里的前后两环。

前一节 \lstinline|0-effective-mass| 做的事情很简单：从 correlator 直接看出 plateau 在哪里、噪声什么时候开始压过信号、bootstrap 误差条是否和样本波动一致。后一节 \lstinline|1-c2pt-fit| 做的事情更进一步：在你已经看清楚 plateau 的前提下，用一个显式模型去提取 ground state energy 和 overlap。

这两步之间的逻辑关系是：

\begin{enumerate}[leftmargin=2em]
  \item \lstinline|0-effective-mass| 先告诉你“哪一段时间窗看起来像单态区”；
  \item \lstinline|1-c2pt-fit| 再把这段时间窗写成具体的拟合区间；
  \item 如果 \lstinline|0-effective-mass| 根本没有 plateau，\lstinline|1-c2pt-fit| 里的任何漂亮数字都不可信；
  \item 如果 \lstinline|0-effective-mass| 的后端误差已经炸开，那么 \lstinline|1-c2pt-fit| 不应该把那一段硬塞进窗口里。
\end{enumerate}

因此，\lstinline|0-effective-mass| 不是“前菜”，而是 \lstinline|1-c2pt-fit| 的诊断图。它告诉学生：你接下来应该在哪些时间点上认真建模，哪些时间点只适合用来说明“不该拟合到这里”。

\paragraph{推荐的阅读和动手顺序}

\begin{enumerate}[leftmargin=2em]
  \item 先回看 \lstinline|0-effective-mass| 的图，标出你认为最稳定的 plateau 区域；
  \item 再打开 \lstinline|1-c2pt-fit| 的 notebook，看它是如何把这个窗口变成 \lstinline|tmin/tmax|；
  \item 对比不同 \lstinline|tmin| 下的 \lstinline|chi^2/\mathrm{dof}|、能量和 overlap；
  \item 把 notebook 里试出来的稳定窗口写进 input file；
  \item 重新跑一遍，检查 notebook 和 input file 是否给出一致结论。
\end{enumerate}

\paragraph{带着问题去练习}

学生在从 \lstinline|0-effective-mass| 过渡到 \lstinline|1-c2pt-fit| 时，至少要能回答下面这些问题：

\begin{enumerate}[leftmargin=2em]
  \item 为什么 plateau 区域比前端更适合做 fit？
  \item 为什么后端即使“看起来更平”，也不一定更可信？
  \item 如果 plateau 只有很短一段，应该优先缩窄窗口还是增加 state 数？
  \item 为什么 k0 和 k6 的最佳窗口可能不同？
  \item 为什么 notebook 中试出来的窗口必须被 input file 冻结下来，才能算真正复现？
\end{enumerate}

\paragraph{这一段的练习题}

\begin{enumerate}[leftmargin=2em]
  \item 任选一张 \lstinline|0-effective-mass| 图，圈出你认为最适合进入 \lstinline|1-c2pt-fit| 的候选时间窗，并说明理由。
  \item 任选另一张图，故意把窗口往前移动一个时间点，再比较 \lstinline|chi^2/\mathrm{dof}| 和能量是否变差。
  \item 用一句话说明：为什么 effective mass 只能告诉你“哪里像 plateau”，却不能替代真正的 correlator fit？
  \item 写下你最终选定的 \lstinline|tmin|，并说明这个选择是从图像直觉、统计稳定性还是物理可解释性出发的。
\end{enumerate}

这组练习的目的不是把学生训练成“会报一个数字的人”，而是训练他们先看图、再建模、最后才接受拟合结果。

\subsubsection{0-effective-mass：把概念变成图}

在真正进入两点拟合之前，学生应该先把 \lstinline|examples/l64c64a076_m140/analysis/0-effective-mass/| 跑通。这个阶段只做一件事：从两点 correlator 直接构造 effective mass，并把结果画出来。它不是最终物理提取，而是把前面学到的统计概念落到一张图上。

\paragraph{运行方式}

这个最小练习有两种等价的跑法：

\begin{itemize}[leftmargin=2em]
  \item \textbf{notebook 方式}：打开 \lstinline|examples/l64c64a076_m140/analysis/0-effective-mass/meff_k0.ipynb| 和 \lstinline|meff_k6.ipynb|，逐格执行，先看输入数据、再看 effective mass 公式、最后看图。
  \item \textbf{input file 方式}：打开 \lstinline|examples/l64c64a076_m140/input_files/0-effective-mass/meff_k0.txt| 和 \lstinline|meff_k6.txt|，用同一套底层 CLI 跑一遍，把 notebook 里调出来的参数冻结成可复现版本。
\end{itemize}

两种方式的分工很明确：notebook 负责理解，input file 负责复现。前者允许学生临时检查中间量，后者要求参数固定、输出稳定。

\paragraph{预期结果}

跑完之后，学生应该看到三类结果：

\begin{itemize}[leftmargin=2em]
  \item 一张或两张 effective mass 图，图上能看出前端弯曲、中间 plateau、后端噪声放大这三种行为；
  \item bootstrap 误差条已经正确反映样本波动，而不是只剩一条平线或完全失控；
  \item k0 和 k6 的图在 plateau 起点、噪声增长速度或误差大小上可能不同，这正是后续选窗口要利用的信息。
\end{itemize}

如果跑完以后图像没有 plateau、误差条明显不合理，或者 notebook 和 input file 的结果对不上，那就说明这一步还没有学透，不应该急着进入 \lstinline|1-c2pt-fit|。

这一步最适合训练学生的地方，是让他们在图像里看到三个统计事实：

\begin{itemize}[leftmargin=2em]
  \item 小 $t$ 处 excited-state contamination 往往更强；
  \item 中间可能出现一段近似平坦的 plateau；
  \item 大 $t$ 处误差条快速变大，读到的往往主要是噪声。
\end{itemize}

Week 1 的训练目标有四个，而且每一个都要对应一次马上动手的练习：

\begin{enumerate}[leftmargin=2em]
  \item 识别 correlator 在不同 $t$ 区间的噪声增长；
  \item 通过 effective mass 观察 plateau 是否开始形成；
  \item 理解 bootstrap 平均值、误差条和样本波动之间的关系；
  \item 学会从图上判断哪些时间点适合以后进入 \lstinline|1-c2pt-fit| 的拟合窗口。
\end{enumerate}

这个练习的核心不是“算一个漂亮的数”，而是让学生先回答下面这些问题：

\begin{itemize}[leftmargin=2em]
  \item effective mass 在小 $t$ 和大 $t$ 的行为为什么不同；
  \item 哪些点看起来已经接近 plateau；
  \item 误差条什么时候开始快速变大；
  \item k0 和 k6 的有效质量图在噪声结构上有什么不同；
  \item 为什么 effective mass 只能作为窗口选择的诊断，而不是最终物理结果。
\end{itemize}

对学生来说，\lstinline|0-effective-mass| 的正确用法是：

\begin{enumerate}[leftmargin=2em]
  \item 先跑 notebook 版，目的是看中间量、熟悉图形、观察 plateau；
  \item 再跑 input file 版，确认同样的数据、同样的参数会得到同样的 effective-mass 结果；
  \item 把这一步的观察结果记下来，作为后面 \lstinline|1-c2pt-fit| 选择拟合窗口的依据。
\end{enumerate}

如果学生连 effective mass 图都看不懂，就不要急着进入两点拟合，更不要直接跳到 DA ratio。

\subsubsection{如何读 effective mass 图}

一张 effective mass 图通常要同时看三件事：均值、误差条和时间依赖趋势。学生第一次看图时，不要急着问“这个数是不是对”，先问下面几个更基本的问题：

\begin{enumerate}[leftmargin=2em]
  \item 图的前端是不是还在明显弯曲？如果是，说明 excited-state contamination 还比较强；
  \item 中间有没有一段比较平的区域？如果有，这往往就是后面 \lstinline|1-c2pt-fit| 的候选窗口；
  \item 图的后端误差是不是迅速放大？如果是，说明继续往后看已经主要是在读噪声；
  \item 不同 momentum 的图是不是有不同的 plateau 起点？如果有，说明拟合窗口不能机械统一；
  \item bootstrap 误差是不是和肉眼看到的波动一致？如果不一致，要先怀疑统计处理，而不是急着解释物理。
\end{enumerate}

effective mass 图最重要的用途不是“报一个最终值”，而是帮助学生学会区分三种区域：

\begin{itemize}[leftmargin=2em]
  \item \textbf{前端}：信号仍然受激发态污染，数值变化快，但物理解释不稳；
  \item \textbf{平台区}：局部斜率近似稳定，是后续拟合最有价值的区域；
  \item \textbf{后端}：噪声快速上升，图像看似平，但那种“平”往往只是统计涨落主导。
\end{itemize}

学生在 Week 1 中要练的，就是把这三类区域分辨出来。

\subsubsection{Week 1 的练习与检查}

为了确保学生不是“看过一遍就算学了”，Week 1 的前半段应该至少完成下面的练习：

\begin{enumerate}[leftmargin=2em]
  \item 分别画出 k0 和 k6 的 effective mass 图，并用一句话描述它们 plateau 的差别；
  \item 标出每张图中你认为最可能进入 \lstinline|1-c2pt-fit| 的时间窗；
  \item 观察 bootstrap 误差条在不同 $t$ 的变化，解释为什么后端更难拟合；
  \item 对 notebook 和 input file 两种方式各跑一次，确认输出图和表格一致；
  \item 写一段简短记录，说明你为什么认为某个窗口更适合拟合，而不是另一个窗口。
\end{enumerate}

这些练习的目的不是让学生“猜对标准答案”，而是训练他们形成一致的判断逻辑：看到图 -> 找 plateau -> 想噪声 -> 再联系拟合窗口。

\subsubsection{DA ratio 的最小物理图像}

Month 1 的学生还需要知道后面的 DA 流程是如何接在两点拟合之后的。当前代码里的最小 ratio 结构是

\begin{equation}
R(t) = \frac{\mathcal{N}(t)}{C_{2\mathrm{pt}}(t)}.
\end{equation}

对于 \lstinline|T5| 通道，分子有最小形式

\begin{equation}
\mathcal{N}_{T5}(t) \propto
\sum_n \frac{1}{2}\sqrt{A_n 2E_n}\,m_n
\left(e^{-E_n t} - e^{-E_n (N_t-t)}\right),
\end{equation}

而 \lstinline|Z5| 通道额外带有动量因子

\begin{equation}
 \mathcal{N}_{Z5}(t) \propto
\sum_n \frac{1}{2}\sqrt{A_n 2E_n}\,\frac{P_z}{E_n}\,m_n
\left(e^{-E_n t} - e^{-E_n (N_t-t)}\right).
\end{equation}

这就是为什么 Month 1 必须先把两点拟合学透：DA ratio 的模型参数直接依赖两点拟合得到的 $A_n$ 和 $E_n$。

\paragraph{这一步到底在做什么}

\lstinline|1-c2pt-fit| 不是在“把一条曲线硬塞给数据”，而是在一个已经被 \lstinline|0-effective-mass| 诊断过的时间窗里，回答两个具体问题：

\begin{enumerate}[leftmargin=2em]
  \item ground-state energy 是多少；
  \item ground-state overlap 是多少。
\end{enumerate}

换句话说，\lstinline|0-effective-mass| 负责告诉你“哪里像单态区”，而 \lstinline|1-c2pt-fit| 负责把“像单态区”这件事变成可复现的参数提取。前者给直觉，后者给数值；前者看图，后者建模。

因此，学生在进入 \lstinline|1-c2pt-fit| 之前，至少应该已经能从 effective mass 图上指出：

\begin{itemize}[leftmargin=2em]
  \item 哪一段时间窗最像 plateau；
  \item 为什么更早的点不适合直接拿来拟合；
  \item 为什么更晚的点虽然噪声可能更大，但不一定更可信；
  \item 为什么不同 momentum 往往需要不同的窗口。
\end{itemize}

\paragraph{这一节的练习题}

\begin{enumerate}[leftmargin=2em]
  \item 先看 \lstinline|0-effective-mass| 的图，给 k0 和 k6 各选一个你认为最稳的候选窗口，并写下理由。
  \item 只允许你改一个参数：如果把 \lstinline|tmin| 向前或向后移动一个时间点，你预期 \lstinline|chi^2/\mathrm{dof}| 会怎么变，为什么？
  \item 比较 \lstinline|1-state| 和 \lstinline|2-state| 的任务差别：前者为什么更适合先做窗口扫描，后者为什么更依赖 prior？
  \item 如果 \lstinline|0-effective-mass| 显示 plateau 很短，你会先考虑扩大窗口还是增加 state 数？把你的判断逻辑写出来。
\end{enumerate}

这组练习的目标，是让学生带着问题进入拟合，而不是先把拟合跑完再回头解释结果。

\subsubsection{1-c2pt-fit：从 plateau 到参数提取}

\lstinline|1-c2pt-fit| 是 Week 1 里真正把“看图”变成“提参”的那一步。它不是重复 \lstinline|0-effective-mass|，而是用一个明确的 correlator 模型，把已经看起来接近 plateau 的那段区间转成 ground-state energy 和 overlap 的数值估计。

对两点函数来说，最基本的物理图像是：在有限时间延拓下，correlator 可以写成若干个态的叠加。最简的单态近似当然是

\begin{equation}
C_{2\mathrm{pt}}(t) \approx A_0\,e^{-E_0 t},
\end{equation}

但真实数据在早时刻往往还混有 excited states，所以更现实的模型是

\begin{equation}
C_{2\mathrm{pt}}(t) = \sum_{n=0}^{N-1} A_n\,K_n(t),
\end{equation}

其中 kernel \(K_n(t)\) 的具体形式取决于 correlator 类型。对当前 example 来说，学生要先记住的不是每个符号的代码名字，而是下面这件事：\textbf{更高 state 数并不是“越多越好”，而是“是否足以解释你选定的时间窗”}。

因此，\lstinline|1-c2pt-fit| 的任务可以拆成三个层次：

\begin{enumerate}[leftmargin=2em]
  \item 先选一个看起来已经接近单态区的窗口；
  \item 再用 1-state / 2-state / 3-state 模型去解释这段窗口；
  \item 最后比较不同模型下的参数稳定性、\lstinline|chi^2/\mathrm{dof}| 和 bootstrap 波动。
\end{enumerate}

这也是为什么 \lstinline|1-c2pt-fit| 的学习顺序必须跟着 \lstinline|0-effective-mass| 走：前者不先看图，后者就没有合理的窗口起点。

\paragraph{运行方式}

\begin{itemize}[leftmargin=2em]
  \item \textbf{notebook 方式}：打开 \lstinline|examples/l64c64a076_m140/analysis/1-c2pt-fit/| 里的 notebook，先观察 effective mass，再看拟合窗口、state 数和输出表格。
  \item \textbf{input file 方式}：打开对应的 \lstinline|examples/l64c64a076_m140/input_files/| 下的两点拟合输入文件，把 notebook 里试出来的窗口和 prior 固化下来。
\end{itemize}

Notebook 的角色是“试”，input file 的角色是“定”。学生应该先在 notebook 里看懂不同窗口为什么会改变结果，再把最终选择写进 input file，确认结果能被重复得到。

\paragraph{预期结果}

跑完 \lstinline|1-c2pt-fit| 后，学生应该至少看到以下几类输出：

\begin{itemize}[leftmargin=2em]
  \item 一张两点 correlator 的拟合图，图上能看到数据点、拟合曲线和候选窗口；
  \item 一个结果表，里面有 ground-state energy、overlap、\lstinline|chi^2/\mathrm{dof}| 和 p-value；
  \item 不同 state 数或不同 \lstinline|tmin| 下结果的对比；
  \item bootstrap 汇总后的中心值和不确定度。
\end{itemize}

如果图上没有明显的拟合区间，或者不同输入方式得到的结果不一致，那么这一步还没有真正学会，不应该急着进入 DA ratio。

\paragraph{这一节的练习题}

\begin{enumerate}[leftmargin=2em]
  \item 先从 \lstinline|0-effective-mass| 的图里选出一个你认为最稳的窗口，再在 \lstinline|1-c2pt-fit| 里验证这个窗口是否真的稳定。
  \item 比较 1-state 和 2-state：如果 1-state 的 \lstinline|chi^2/\mathrm{dof}| 很高，而 2-state 明显下降，这说明了什么？
  \item 故意把 \lstinline|tmin| 改到更早一点，观察结果如何变化，并解释这通常意味着什么物理污染。
  \item 把 notebook 里调过的参数写进 input file，再跑一遍；如果输出不同，先检查哪里没有冻结。
\end{enumerate}

\paragraph{这一步应该形成的判断习惯}

学生在 \lstinline|1-c2pt-fit| 中要养成的习惯是：

\begin{itemize}[leftmargin=2em]
  \item 先看 plateau 再看拟合值；
  \item 先看窗口再看 \lstinline|chi^2/\mathrm{dof}|；
  \item 先比较不同 state 数，再决定是否加复杂度；
  \item 先把 notebook 理顺，再把结果冻结到 input file。
\end{itemize}

\paragraph{为什么窗口、state 数和 prior 必须一起看}

这三件事不是三个互不相干的按钮，而是一组彼此约束的统计选择。

\begin{itemize}[leftmargin=2em]
  \item \textbf{窗口} 决定你把哪些时间点当成“可被模型解释”的数据；
  \item \textbf{state 数} 决定你愿意用多复杂的物理模型去解释这段数据；
  \item \textbf{prior} 决定你在数据不够强时，愿意借用多少外部物理知识来稳定拟合。
\end{itemize}

这三者的关系可以理解成一个三角形：

\begin{itemize}[leftmargin=2em]
  \item 窗口太短，拟合可能很稳，但参数会偏；
  \item 窗口太长，后端噪声会把拟合拖垮；
  \item state 数太少，模型解释不了 plateau 前后的弯曲；
  \item state 数太多，参数会变得不稳定，甚至把噪声拟合成物理；
  \item prior 太弱，拟合会漂；prior 太强，又可能把数据“压”回先验中心。
\end{itemize}

所以学生不要把这些参数看成独立可调项。正确的顺序是：先用 \lstinline|0-effective-mass| 判断大概的窗口，再用 \lstinline|1-state| 看地基是否稳，再考虑 \lstinline|2-state| 或 \lstinline|3-state|，最后才决定需不需要 prior 来固定或约束低态。

在本项目里，\lstinline|low_state_prior_tmin[pz]| 的角色尤其重要：它不是“随便挑一行表”，而是告诉下一层拟合“上层里哪一个稳定结果最值得拿来当先验”。同理，\lstinline|pz0_ground_energy| 也不是装饰参数，而是给 ground state 一个物理上合理的起点，避免优化器在明显不合理的区域里乱跑。

\paragraph{这一节的练习题}

\begin{enumerate}[leftmargin=2em]
  \item 如果你把窗口往前推一点，\lstinline|chi^2/\mathrm{dof}| 可能变小也可能变大。请用 excited-state contamination 的语言解释这两种情况各意味着什么。
  \item 如果你把窗口往后推一点，误差条通常会变大。请解释为什么这并不自动意味着结果更好。
  \item 比较 1-state 和 2-state：在什么情况下“多一个 state”是真正改善模型，在什么情况下只是让参数更多？
  \item 为什么 prior 只能帮助稳定拟合，而不能替代窗口选择？请给出你的物理理由，而不是只写“因为代码这么做”。
\end{enumerate}

\subsubsection{两点拟合的窗口、状态数和先验选择}

对学生来说，最重要的不是记住某一组固定数字，而是理解当前两点拟合的决策顺序。

\paragraph{推荐顺序}
\begin{enumerate}[leftmargin=2em]
  \item 先跑 \lstinline|1-state|，只看 ground-state 是否稳定；
  \item 在 \lstinline|1-state| 的 plateau 上选一个稳定的 \lstinline|tmin|；
  \item 用这个稳定行作为 \lstinline|low_state_prior_tmin[pz]| 的来源；
  \item 再跑 \lstinline|2-state|；
  \item 只有当 \lstinline|2-state| 不能解释数据时，才考虑 \lstinline|3-state|。
\end{enumerate}

\paragraph{窗口参数}
\begin{itemize}[leftmargin=2em]
  \item \lstinline|tmin_window| 决定扫描哪些起点，主要用来避开早时刻的 excited-state contamination；
  \item \lstinline|tmax| 决定晚时刻是否已经进入噪声主导区；
  \item 在实际使用中，\lstinline|tmin| 通常比 \lstinline|tmax| 更重要；
  \item 更高 state 数的窗口通常比低 state 数更窄，因为模型更容易过拟合。
\end{itemize}

\paragraph{先验和锚定}
当前实现把 two-point 的关键物理约束拆成两类：

\begin{itemize}[leftmargin=2em]
  \item \lstinline|pz0_ground_energy|：来自可信的 $p_z=0$ 参考值，用来给 $E_0$ 提供初始猜测；如果打开 \lstinline|fix_ground_energy_from_dispersion|，它还会直接固定 $E_0$；
  \item \lstinline|low_state_prior_tmin[pz]|：从前一个 state 的 fit table 中挑选某一行，作为下一个 state 的 prior 来源。
\end{itemize}

这两个量的角色不同：

\begin{itemize}[leftmargin=2em]
  \item \lstinline|pz0_ground_energy| 是物理锚点，回答“ground state 大概应该在哪”；
  \item \lstinline|low_state_prior_tmin| 是表格行选择器，回答“上一层拟合里哪一行最适合作为下一层的先验”。
\end{itemize}

\paragraph{状态数如何递增}
\begin{itemize}[leftmargin=2em]
  \item \lstinline|1-state|：没有低态先验，只做最基础的 ground-state 提取；
  \item \lstinline|2-state|：通常对 $E_0$ 加 prior 或固定锚定，并让 $E_1$ 吸收 excited-state correction；
  \item \lstinline|3-state|：在 \lstinline|2-state| 的基础上再多一个 excited level，先验的作用更强，窗口也应该更保守。
\end{itemize}

\paragraph{协方差与 shrinkage}
当前代码优先尝试 correlated fit。若协方差矩阵不稳定，代码会按照
\begin{equation}
\lambda \in \{0.0, 0.1, 0.2, 0.3, 0.5, 1.0\}
\end{equation}
的顺序尝试 shrinkage。这里：

\begin{itemize}[leftmargin=2em]
  \item $\lambda=0$ 表示完全保留协方差；
  \item $\lambda=1$ 表示只保留对角项；
  \item 中间值是在稳定性和相关性信息之间做折中。
\end{itemize}

学生应该理解：correlated fit 更接近真实统计结构，但也更依赖稳定协方差；如果协方差不稳，先用 shrinkage 稳住数值，再讨论物理解释。

\subsubsection{DA ratio 的窗口和诊断逻辑}

DA ratio 的拟合不是“随便拿一段时间窗做一个比值”，而是一个和 two-point 拟合强耦合的流程。

\paragraph{当前输入约束}
\begin{itemize}[leftmargin=2em]
  \item \lstinline|fit_target = ratio|；
  \item \lstinline|fit_component| 可以是 \lstinline|real|、\lstinline|imag| 或 \lstinline|both|；
  \item 真实 example 只保留 \lstinline|b_X|、\lstinline|OT5|、\lstinline|bT = 0|；
  \item \lstinline|decay_constant_check = true| 时，代码会强制只看 real part，并把 \lstinline|bTlist| 与 \lstinline|bzlist| 收缩到 \lstinline|0|；
  \item \lstinline|two_point_fit_root| 和 \lstinline|two_point_fit_window_by_pz| 决定 DA ratio 使用哪一行 two-point 参考。
\end{itemize}

\paragraph{窗口选择}
DA 的 \lstinline|fit_window| 是核心调节钮。正常模式下，学生应该直接指定每个 momentum 的基准窗口；\lstinline|decay_constant_check| 模式下，代码会围绕基准窗口做一个小范围扫描，用来检查：

\begin{itemize}[leftmargin=2em]
  \item 抽出的 decay constant 是否对 \lstinline|tmin| 稳定；
  \item \lstinline|tmax| 是否已经进入噪声主导区；
  \item 结果是否明显受 excited-state contamination 影响；
  \item 不同窗口之间的 \lstinline|chi^2/dof| 和相对误差是否相互一致。
\end{itemize}

\paragraph{两点与 DA 的连接}
DA ratio 不是独立于 two-point 的。当前代码里，two-point 的拟合结果直接进入 DA ratio 的模型：

\begin{itemize}[leftmargin=2em]
  \item two-point 的 $A_n$ 和 $E_n$ 决定 DA 分子和分母的时间依赖；
  \item \lstinline|T5| 和 \lstinline|Z5| 只是不同的算符通道，底层都依赖同一套 two-point 约束；
  \item \lstinline|two_point_fit_sample_coupled| 若打开，则每个 bootstrap 样本会尽量使用对应的 two-point sample 行，而不是只看 summary table。
\end{itemize}

所以 Week 1 的思路应该是：先把两点拟合窗口选稳，再让 DA ratio 在这个稳态上做提取，而不是反过来。

\subsubsection{Week 1 和后续 example 的连接}

当前仓库中的实际分析顺序是：

\begin{center}
\begin{tabular}{p{0.18\linewidth}p{0.30\linewidth}p{0.42\linewidth}}
\toprule
阶段 & example 目录 & 核心任务 \\
\midrule
0-effective-mass & \lstinline|examples/.../analysis/0-effective-mass/| & 先只算 effective mass 和画图，练习 plateau、bootstrap 和噪声诊断 \\
1-c2pt-fit & \lstinline|examples/.../analysis/1-c2pt-fit/| & 提取两点能量和 overlap，建立窗口选择规则 \\
2-bm & \lstinline|examples/.../analysis/2-bm/| & 提取 DA bare matrix element \\
3-normalize & \lstinline|examples/.../analysis/3-normalize/| & 归一化，当前只保留 mode3 \\
4-FT & \lstinline|examples/.../analysis/4-FT/| & 做 Fourier 变换，检查最终分布 \\
\bottomrule
\end{tabular}
\end{center}

Week 1 只要求学生真正吃透第一阶段，但应该知道后面三步为什么存在、每一步为什么都依赖前一步。

\subsection{如何运行当前 example：先看原则，细节放到后面的实操部分}

本仓库当前保留了两种等价但用途不同的运行方式：notebook 和 input file。

这里先给学生一个总原则：

\begin{itemize}[leftmargin=2em]
  \item \textbf{先 notebook，后 input file}；
  \item notebook 用来理解流程、检查中间量、调窗口；
  \item input file 用来冻结参数、复现结果、提交最终版本；
  \item 两者必须走同一条底层分析链，否则结果不应被认为是同一个 example。
\end{itemize}

具体到当前 example，真实的操作顺序和命令细节放在下面的实操任务里。

\subsection{Week 1 的实操主线：环境、代码理解与 golden example 复现}

\subsubsection{Week 1 目标}

第一周的目标是\textbf{复现已有分析流程，并把流程看懂、记录清楚}。这周不要急着重构，不要急着做 agent。先要回答三个问题：

\begin{enumerate}[leftmargin=2em]
  \item 这个 meson DA 分析样例从哪些输入开始？
  \item 中间经过哪些步骤？
  \item 最终产生哪些物理结果、图和检查量？
\end{enumerate}

\subsubsection{Week 1 的学习-练习循环}

Week 1 不应该被理解成“先把工具书看完，再统一上手”。更合理的做法是：每学一个分析模块，就立刻把它放进 notebook 或 input file 里验证一次，再把验证结果写回文档。

学生需要边学边练的内容包括：

\begin{itemize}[leftmargin=2em]
  \item Git 基本操作；
  \item Python 虚拟环境或 conda 环境；
  \item 项目目录结构；
  \item 如何运行现有 notebook 或 script，并立刻核对输出是否符合预期；
  \item meson DA 分析中的基本对象：2pt、3pt/ratio、bare matrix element、$z$ dependence、$P_z$ dependence、real/imag symmetry；
  \item 什么是 golden example，为什么它可以作为 regression case；
  \item 如何把 notebook 里摸索出来的参数冻结到 input file 里。
\end{itemize}

\subsubsection{Week 1 实操主线：先 notebook，再 input file}

Week 1 的实操不是“随便点一个 notebook 看看”，而是按固定顺序把同一条分析链跑两遍。

\paragraph{第一遍：Notebook 探索}

第一遍的目标是看懂流程和中间量，不追求最后一步就完全冻结成最终版。

\begin{enumerate}[leftmargin=2em]
  \item 先打开 \lstinline|examples/l64c64a076_m140/analysis/0-effective-mass/| 里的 notebook；
  \item 先看 effective mass 的 bootstrap 结果和图，熟悉 plateau、噪声和误差条；
  \item 再打开 \lstinline|examples/l64c64a076_m140/analysis/1-c2pt-fit/| 里的 notebook；
  \item 从 \lstinline|1-c2pt-fit| 开始看两点 correlator 的 plateau 和拟合窗口；
  \item 选出稳定的 \lstinline|tmin|，确认 \lstinline|chi^2/\mathrm{dof}|、能量和 overlap 没有明显异常；
  \item 再打开 \lstinline|examples/l64c64a076_m140/analysis/2-bm/|，用同一套 two-point 参考去看 DA bare matrix element；
  \item 接着打开 \lstinline|examples/l64c64a076_m140/analysis/3-normalize/|，只保留 \lstinline|mode3| 路径；
  \item 最后打开 \lstinline|examples/l64c64a076_m140/analysis/4-FT/|，检查最终 Fourier 结果的形状和物理趋势。
\end{enumerate}

Notebook 这一遍的重点是：

\begin{itemize}[leftmargin=2em]
  \item 理解每一步输出什么；
  \item 找到每一步里最敏感的参数；
  \item 把 notebook 里临时调过的窗口记下来；
  \item 不要急着改代码结构。
\end{itemize}

\paragraph{第二遍：Input file 冻结与复现}

第二遍的目标是把 notebook 里已经理解清楚的参数，写进 input file，然后用 CLI 重跑，确认结果一致。

当前 example 里常见的 input file 是：

\begin{itemize}[leftmargin=2em]
  \item \lstinline|examples/l64c64a076_m140/input_files/0-effective-mass/meff_k0.txt|；
  \item \lstinline|examples/l64c64a076_m140/input_files/0-effective-mass/meff_k6.txt|；
  \item \lstinline|examples/l64c64a076_m140/input_files/2-bm/da_nstate_k0.txt|；
  \item \lstinline|examples/l64c64a076_m140/input_files/2-bm/da_nstate_k6.txt|；
  \item \lstinline|examples/l64c64a076_m140/input_files/3-normalize/da_normalize_mode3.txt|；
  \item \lstinline|examples/l64c64a076_m140/input_files/4-FT/da_fourier_mode3.txt|。
\end{itemize}

Input file 的运行要求是：

\begin{enumerate}[leftmargin=2em]
  \item 参数必须和 notebook 一致，尤其是数据路径、窗口、state 数、模式和输出目录；
  \item 运行顺序必须和 notebook 的分析链一致，先两点，再 bare matrix element，再 normalize，最后 Fourier；
  \item 跑完以后要检查输出文件名、表格和图是否与 notebook 对齐；
  \item 如果不一致，先检查 input file 是否漏了参数，再检查 notebook 是否靠了手动状态。
\end{enumerate}

对学生最重要的判断标准是：\textbf{notebook 和 input file 应该得到同一物理结论}。如果二者不一致，就说明参数冻结还不彻底，不能把结果当作 final version。

\subsubsection{Week 1 每日任务单}

Week 1 的后半段采用“先看当天目标，再完成当天练习，再交当天产出”的方式推进。每一天都尽量只抓一件主事：先把当天该学的内容吃透，再把它变成一个可检查、可复现的小结果。

\begin{longtable}{p{0.11\textwidth}p{0.28\textwidth}p{0.29\textwidth}p{0.26\textwidth}}
\toprule
天数 & 当天主题 & 当天主要产出 & 当天检查点 \\
\midrule
Day 1 & 环境准备与 repo 初读 & \texttt{docs/repo\_initial\_notes.md} & 知道 repo 结构和 golden example 在哪里 \\
Day 2 & 手动复现 golden example & \texttt{docs/golden\_example\_reproduction.md} 与运行 log & 能从干净环境复现一次主流程 \\
Day 3 & 画 data-flow map & \texttt{docs/data\_flow\_map.md} & 能把流程拆成输入、输出和物理含义 \\
Day 4 & 识别 config 参数 & \texttt{docs/config\_design\_notes.md} 和 draft YAML & 能指出哪些值必须 config 化 \\
Day 5 & 总结与 review & \texttt{docs/week1\_report.md} & 能说明自己做了什么、还卡在哪里 \\
\bottomrule
\end{longtable}

\paragraph{Day 1：环境准备与 repo 初读}

课堂动作：

\begin{enumerate}[leftmargin=2em]
  \item clone 或进入项目 repo；
  \item 创建 Python 环境；
  \item 安装依赖；
  \item 运行现有测试；
  \item 阅读 README 和 examples 目录；
  \item 写下当前目录结构说明。
\end{enumerate}

建议命令：

\begin{lstlisting}[language=bash]
git status
python --version
pip list
pytest
find . -maxdepth 3 -type f | sort | head -100
\end{lstlisting}

要交的东西：

\begin{lstlisting}
docs/repo_initial_notes.md
\end{lstlisting}

内容包括：

\begin{itemize}[leftmargin=2em]
  \item 当前 repo 里有哪些主要目录；
  \item 哪些目录是数据、代码、配置、输出、文档；
  \item 当前是否能成功运行测试；
  \item 当前不理解的问题列表。
\end{itemize}

Codex prompt 示例：

\begin{lstlisting}
We are starting an AI-assisted LQCD meson DA analysis-agent project.
Please inspect the repository structure and write a short developer-oriented summary in `docs/repo_initial_notes.md`.

Focus on:
- main code directories;
- example analysis directories;
- likely inputs and outputs;
- existing scripts or notebooks;
- tests if any.

Do not modify source code. Only create or update the documentation file.
\end{lstlisting}

验收：

\begin{itemize}[leftmargin=2em]
  \item 学生能说清楚 repo 大致结构；
  \item 知道 golden example 在哪里；
  \item 知道如何运行最基本命令；
  \item 有第一份文档记录。
\end{itemize}

\paragraph{Day 2：手动复现 golden example}

课堂动作：

\begin{enumerate}[leftmargin=2em]
  \item 按现有方式运行 golden example；
  \item 记录每一条命令；
  \item 记录需要哪些输入文件；
  \item 记录生成哪些输出文件；
  \item 保存运行 log。
\end{enumerate}

要交的东西：

\begin{lstlisting}
docs/golden_example_reproduction.md
outputs/golden_example_manual/run.log
\end{lstlisting}

文档中必须包含：

\begin{itemize}[leftmargin=2em]
  \item 运行日期；
  \item git commit hash；
  \item Python 环境；
  \item 输入文件路径；
  \item 运行命令；
  \item 输出文件列表；
  \item 运行中遇到的问题。
\end{itemize}

建议记录 git commit：

\begin{lstlisting}[language=bash]
git rev-parse HEAD
\end{lstlisting}

Codex prompt 示例：

\begin{lstlisting}
I manually ran the golden example. Below are the commands and output files.
Please help me turn my notes into a clean reproduction document at `docs/golden_example_reproduction.md`.

Requirements:
- include environment information;
- include exact commands;
- include input files;
- include output files;
- include unresolved questions;
- do not modify code.

Here are my raw notes:
[paste notes]
\end{lstlisting}

验收：

\begin{itemize}[leftmargin=2em]
  \item 能从干净环境复现一次 golden example；
  \item 所有命令被记录；
  \item 输出文件被列出；
  \item 失败或 warning 没有被隐藏。
\end{itemize}

\paragraph{Day 3：画 data-flow map}

课堂动作：

把 golden example 拆成流程图或表格，例如：

\begin{lstlisting}
raw data
  -> data loading
  -> 2pt fit input
  -> ratio construction
  -> ratio fit
  -> bare matrix element
  -> plots and summary
\end{lstlisting}

要交的东西：

\begin{lstlisting}
docs/data_flow_map.md
\end{lstlisting}

建议表格：

\begin{longtable}{p{0.18\textwidth}p{0.28\textwidth}p{0.28\textwidth}p{0.16\textwidth}}
\toprule
步骤 & 输入 & 输出 & 物理/技术含义 \\
\midrule
Data audit & raw correlator files & file checklist & 检查数据是否齐全 \\
2pt fit & two-point correlator & mass/energy fit result & 提供能量和 overlap 信息 \\
Ratio analysis & 3pt/2pt ratio & bare matrix element & 提取矩阵元 \\
Plotting & summary table & figures & 检查稳定性和物理趋势 \\
\bottomrule
\end{longtable}

Codex prompt 示例：

\begin{lstlisting}
Please create `docs/data_flow_map.md` for the golden example.

Use the current analysis notes and file list below.
The document should explain:
- each analysis step;
- input files for each step;
- output files for each step;
- what physical or technical quantity each step produces;
- which steps should later become separate pipeline stages.

Do not modify source code.

Notes:
[paste notes]
\end{lstlisting}

验收：

\begin{itemize}[leftmargin=2em]
  \item 学生能画出完整分析流程；
  \item 知道哪些步骤可以 later 变成 CLI module；
  \item 知道哪些输出应该进入 summary.csv/result.json；
  \item 知道哪些检查属于 physics sanity check。
\end{itemize}

\paragraph{Day 4：识别 config 参数和硬编码路径}

课堂动作：

检查现有脚本或 notebook 中哪些值应该变成 config。

典型 config 字段包括：

\begin{itemize}[leftmargin=2em]
  \item ensemble name；
  \item input data root；
  \item output directory；
  \item momentum list；
  \item $z$ range；
  \item bootstrap/jackknife 参数；
  \item fit model；
  \item fit window；
  \item plot options；
  \item random seed；
  \item overwrite policy。
\end{itemize}

要交的东西：

\begin{lstlisting}
docs/config_design_notes.md
configs/golden_example_draft.yaml
\end{lstlisting}

Codex prompt 示例：

\begin{lstlisting}
Inspect the golden example scripts/notebooks and identify values that should become configuration fields.

Please create:
1. `docs/config_design_notes.md`
2. `configs/golden_example_draft.yaml`

The config should include paths, ensemble metadata, momentum/z selections, fit settings, bootstrap settings, output settings, and overwrite policy.

Do not refactor the analysis code yet.
Do not change existing outputs.
\end{lstlisting}

验收：

\begin{itemize}[leftmargin=2em]
  \item 学生能指出哪些地方硬编码了路径或参数；
  \item 第一版 YAML config 能表达 golden example；
  \item 学生理解 config-driven pipeline 的意义。
\end{itemize}

\paragraph{Day 5：Week 1 总结与导师 review}

课堂动作：

整理第一周报告。

要交的东西：

\begin{lstlisting}
docs/week1_report.md
\end{lstlisting}

报告模板：

\begin{lstlisting}
# Week 1 Report

## 1. Goal

## 2. What I ran

## 3. Inputs and outputs

## 4. Data-flow map

## 5. Config parameters identified

## 6. Problems encountered

## 7. What I learned

## 8. Questions for advisor

## 9. Plan for Week 2
\end{lstlisting}

导师 review 重点：

\begin{enumerate}[leftmargin=2em]
  \item 学生是否真的跑通过 golden example？
  \item 是否记录了完整命令？
  \item 是否理解主要输入输出？
  \item 是否知道哪些参数应该 config 化？
  \item 是否有不理解的问题被明确列出？
\end{enumerate}

\section{Week 2：Config-driven CLI 与 Pydantic schema}

\subsection{Week 2 目标}

第二周的目标是把 golden example 从“手动运行”推进到“config-driven CLI”。

学生要理解：科研分析不能长期依赖 notebook 里手动改路径、手动运行 cell、手动保存图。真正可复现的分析应该由 config 驱动，并且每次运行都留下结构化结果和日志。

\subsection{Week 2 学习内容}

学生需要学习：

\begin{itemize}[leftmargin=2em]
  \item Python package 基础结构；
  \item CLI 入口；
  \item YAML/JSON config；
  \item Pydantic schema；
  \item result.json 和 summary.csv 的区别；
  \item run manifest 的作用。
\end{itemize}

\subsection{Week 2 任务 1：定义 config schema}

创建：

\begin{lstlisting}
src/lqcd_agent/schema/analysis_config.py
\end{lstlisting}

建议 schema 包含：

\begin{lstlisting}[language=Python]
from pydantic import BaseModel, Field
from typing import List, Optional

class EnsembleConfig(BaseModel):
    name: str
    lattice_size: Optional[str] = None
    pion_mass: Optional[float] = None
    notes: Optional[str] = None

class IOConfig(BaseModel):
    input_root: str
    output_dir: str
    allow_overwrite: bool = False

class AnalysisSelection(BaseModel):
    momenta: List[str] = Field(default_factory=list)
    z_values: List[int] = Field(default_factory=list)

class FitConfig(BaseModel):
    model: str
    tmin: Optional[int] = None
    tmax: Optional[int] = None
    bootstrap_samples: int = 500

class GoldenExampleConfig(BaseModel):
    run_id: str
    ensemble: EnsembleConfig
    io: IOConfig
    selection: AnalysisSelection
    fit: FitConfig
\end{lstlisting}

这只是示意，学生应根据真实分析流程修改字段。

Codex prompt 示例：

\begin{lstlisting}
We are converting the golden example into a config-driven CLI.

Please implement Pydantic config schemas in:
`src/lqcd_agent/schema/analysis_config.py`

Use the draft YAML config in `configs/golden_example_draft.yaml` as the reference.

Requirements:
- define nested schemas for ensemble metadata, IO paths, analysis selections, fit settings, and runtime options;
- validate that output_dir is not empty;
- validate that bootstrap_samples is positive;
- validate that z_values are integers;
- add a helper function `load_config(path: str) -> GoldenExampleConfig`;
- add pytest tests for valid config and invalid config.

Do not change the analysis algorithm yet.
\end{lstlisting}

验收：

\begin{itemize}[leftmargin=2em]
  \item YAML config 可以被 load；
  \item 错误 config 会清楚报错；
  \item pytest 覆盖基本 validation；
  \item 学生能解释每个 config 字段的作用。
\end{itemize}

\subsection{Week 2 任务 2：定义 result schema}

创建：

\begin{lstlisting}
src/lqcd_agent/schema/analysis_result.py
\end{lstlisting}

建议 result schema 包含：

\begin{lstlisting}[language=Python]
from pydantic import BaseModel
from typing import Dict, List, Optional

class OutputFiles(BaseModel):
    manifest: str
    result_json: str
    summary_csv: str
    plots: List[str]
    run_log: str

class FitSummary(BaseModel):
    status: str
    model: Optional[str] = None
    chi2_dof: Optional[float] = None
    notes: List[str] = []

class AnalysisResult(BaseModel):
    run_id: str
    status: str
    git_commit: Optional[str]
    config_path: str
    outputs: OutputFiles
    fit_summary: Optional[FitSummary]
    warnings: List[str] = []
    errors: List[str] = []
\end{lstlisting}

Codex prompt 示例：

\begin{lstlisting}
Please implement result schemas for the golden example run in:
`src/lqcd_agent/schema/analysis_result.py`

Requirements:
- include run_id, status, git_commit, config_path;
- include output file paths: manifest, result_json, summary_csv, plots, run_log;
- include fit summary fields if available;
- include warnings and errors;
- add tests that validate a successful result and reject a malformed result.

Do not run the analysis in this task.
\end{lstlisting}

\subsection{Week 2 任务 3：实现最小 CLI}

创建：

\begin{lstlisting}
src/lqcd_agent/cli/run_golden.py
\end{lstlisting}

CLI 的第一版可以只是：

\begin{enumerate}[leftmargin=2em]
  \item 读取 config；
  \item 创建 output directory；
  \item 写 manifest.json；
  \item 调用已有分析脚本或占位函数；
  \item 收集输出文件；
  \item 写 result.json 和 run.log。
\end{enumerate}

建议命令：

\begin{lstlisting}[language=bash]
python -m lqcd_agent.cli.run_golden --config configs/golden_example.yaml
\end{lstlisting}

Codex prompt 示例：

\begin{lstlisting}
Please implement a minimal CLI for running the golden example:
`src/lqcd_agent/cli/run_golden.py`

Requirements:
- command: `python -m lqcd_agent.cli.run_golden --config configs/golden_example.yaml`
- load and validate the Pydantic config;
- create the output directory;
- refuse to overwrite existing output unless allow_overwrite is true;
- write `manifest.json` with run_id, timestamp, git commit, config path, and environment summary;
- write `run.log`;
- for now, call a placeholder function or existing script hook for the actual analysis;
- write `result.json` using the result schema;
- add a CLI smoke test using a temporary directory.

Do not modify raw data.
Do not implement Snakemake yet.
\end{lstlisting}

验收标准：

\begin{itemize}[leftmargin=2em]
  \item CLI 能从 config 启动；
  \item 不允许默认覆盖输出；
  \item manifest.json 存在；
  \item result.json 符合 schema；
  \item CLI smoke test 通过。
\end{itemize}

\subsection{Week 2 任务 4：记录 Codex 使用日志}

学生需要维护：

\begin{lstlisting}
docs/codex_usage_log.md
\end{lstlisting}

记录格式：

\begin{lstlisting}
## Entry 001
Date:
Goal:
Prompt summary:
Files changed:
Tests added:
Did it work? yes/no
What I checked manually:
Problems:
Follow-up prompt:
\end{lstlisting}

目的不是监控学生，而是训练学生学会评估 AI-assisted coding 的质量。

\section{Week 3：数据检查、最小测试体系与输出规范}

\subsection{Week 3 目标}

第三周的目标是让项目从“能运行”变成“更可靠”。具体来说，要加入数据检查、pytest 测试、输出规范和基本 numerical regression 思想。

\subsection{Week 3 学习内容}

学生需要学习：

\begin{itemize}[leftmargin=2em]
  \item data audit 的意义；
  \item pytest 的基本写法；
  \item smoke test、unit test、regression test 的区别；
  \item manifest 和 checksum；
  \item 如何检查输出是否可复现。
\end{itemize}

\subsection{Week 3 任务 1：Data Auditor 模块}

创建：

\begin{lstlisting}
src/lqcd_agent/audit/data_audit.py
\end{lstlisting}

Data Auditor 的职责：

\begin{enumerate}[leftmargin=2em]
  \item 检查 input root 是否存在；
  \item 检查 config 中指定的必要文件是否存在；
  \item 检查输出目录是否安全；
  \item 统计文件数量和基本 metadata；
  \item 写 \texttt{audit\_report.json}；
  \item 不读取或修改大规模 raw data 内容。
\end{enumerate}

Codex prompt 示例：

\begin{lstlisting}
Please implement a Data Auditor module for the golden example.

File:
`src/lqcd_agent/audit/data_audit.py`

Requirements:
- input: validated GoldenExampleConfig;
- check that input_root exists;
- check that required files or directories listed in config exist;
- collect file count and basic metadata if cheap;
- do not modify raw data;
- write `audit_report.json` into output_dir;
- return a structured audit result object;
- add pytest tests using temporary directories.

The auditor should fail clearly if required inputs are missing.
\end{lstlisting}

验收标准：

\begin{itemize}[leftmargin=2em]
  \item 缺文件时能清楚报错；
  \item \texttt{audit\_report.json} 能生成；
  \item 测试覆盖正常和失败情况；
  \item 学生理解为什么 data audit 是 Worker 之前的第一步。
\end{itemize}

\subsection{Week 3 任务 2：最小 pytest 体系}

至少需要以下测试：

\begin{lstlisting}
tests/test_schema.py
tests/test_data_audit.py
tests/test_cli_smoke.py
\end{lstlisting}

测试内容：

\begin{itemize}[leftmargin=2em]
  \item valid config 可以通过；
  \item invalid config 会失败；
  \item data audit 检测缺失输入；
  \item CLI 可以在 temporary output directory 运行；
  \item result.json 可以被 result schema 重新读取。
\end{itemize}

学生需要运行：

\begin{lstlisting}[language=bash]
pytest -q
\end{lstlisting}

Codex prompt 示例：

\begin{lstlisting}
Please add a minimal pytest suite for the Month 1 golden example pipeline.

Tests required:
1. valid config loads successfully;
2. invalid config with negative bootstrap_samples fails;
3. Data Auditor fails when input_root is missing;
4. CLI smoke test writes manifest.json and result.json into a temporary output directory;
5. result.json can be validated by AnalysisResult schema.

Use tmp_path for tests.
Do not require real large raw data in unit tests.
Do not make tests depend on machine-specific paths.
\end{lstlisting}

\subsection{Week 3 任务 3：输出规范}

规定 golden example 每次运行必须输出：

\begin{lstlisting}
manifest.json
result.json
summary.csv
run.log
audit_report.json
plots/
\end{lstlisting}

\subsubsection{manifest.json 应包含}

\begin{itemize}[leftmargin=2em]
  \item \texttt{run\_id}；
  \item timestamp；
  \item git commit；
  \item config path；
  \item output directory；
  \item Python version；
  \item key package versions；
  \item hostname 可选；
  \item random seed 可选。
\end{itemize}

\subsubsection{result.json 应包含}

\begin{itemize}[leftmargin=2em]
  \item run status；
  \item key outputs；
  \item major numerical results；
  \item warnings；
  \item errors；
  \item whether human review is needed later。
\end{itemize}

\subsubsection{summary.csv 应包含}

对于 Month 1，可以先简单记录输出文件和关键数值；后续 Month 2 再扩展为 fit candidate table。

\subsection{Week 3 验收标准}

第三周结束时：

\begin{itemize}[leftmargin=2em]
  \item CLI 已经不是黑盒运行，而是有 manifest/audit/result/log；
  \item pytest 至少有 5 个有效测试；
  \item 学生能解释 smoke test 和 unit test 的区别；
  \item 学生能说明为什么 raw data 不能被 Worker 修改；
  \item 输出目录结构稳定。
\end{itemize}

\section{Week 4：Month 1 集成、最小 Worker 雏形与展示}

\subsection{Week 4 目标}

第四周的目标是完成第一个月的集成展示：学生应该能用一个命令运行 golden example，并生成完整的结构化输出和复现报告。同时，开始引入 Worker Agent 的雏形，但不要求完成完整多 agent 系统。

\subsection{Week 4 学习内容}

学生需要学习：

\begin{itemize}[leftmargin=2em]
  \item Worker Agent 的边界；
  \item task input 和 task output；
  \item 安全规则；
  \item 如何写 demo report；
  \item 如何向导师展示一个科研软件原型。
\end{itemize}

\subsection{Week 4 任务 1：定义 Worker task schema}

创建：

\begin{lstlisting}
src/lqcd_agent/schema/worker_task.py
\end{lstlisting}

示意：

\begin{lstlisting}[language=Python]
from pydantic import BaseModel
from typing import Optional, List

class WorkerTask(BaseModel):
    task_id: str
    task_type: str
    config_path: str
    output_dir: str
    allow_overwrite: bool = False
    notes: Optional[str] = None

class WorkerTaskResult(BaseModel):
    task_id: str
    status: str
    result_json: Optional[str] = None
    summary_csv: Optional[str] = None
    plots: List[str] = []
    warnings: List[str] = []
    errors: List[str] = []
\end{lstlisting}

Codex prompt 示例：

\begin{lstlisting}
Please add a minimal Worker task schema.

File:
`src/lqcd_agent/schema/worker_task.py`

Requirements:
- define WorkerTask with task_id, task_type, config_path, output_dir, allow_overwrite, notes;
- define WorkerTaskResult with status, output paths, warnings, errors;
- validate that task_type is one of the supported values for Month 1: `run_golden_example`;
- add tests for valid and invalid task types.

Do not implement LangGraph yet.
\end{lstlisting}

\subsection{Week 4 任务 2：最小 Worker wrapper}

创建：

\begin{lstlisting}
src/lqcd_agent/worker/run_worker.py
\end{lstlisting}

Worker wrapper 的第一版职责：

\begin{enumerate}[leftmargin=2em]
  \item 读取 WorkerTask；
  \item 检查 task type；
  \item 调用 \texttt{run\_golden} CLI 或内部函数；
  \item 检查输出文件是否存在；
  \item 写 WorkerTaskResult；
  \item 不做 final physics decision。
\end{enumerate}

建议命令：

\begin{lstlisting}[language=bash]
python -m lqcd_agent.worker.run_worker --task tasks/month1_golden_task.yaml
\end{lstlisting}

Codex prompt 示例：

\begin{lstlisting}
Please implement a minimal Worker wrapper for Month 1.

File:
`src/lqcd_agent/worker/run_worker.py`

Requirements:
- command: `python -m lqcd_agent.worker.run_worker --task tasks/month1_golden_task.yaml`;
- load WorkerTask schema;
- support task_type = `run_golden_example`;
- call the existing golden example CLI or internal run function;
- check that result.json, summary.csv, manifest.json, and run.log exist;
- write `worker_result.json`;
- if any required output is missing, mark status as failed and record errors;
- do not decide central physics results;
- do not overwrite existing outputs unless allow_overwrite is true;
- add pytest tests for success and missing config failure.
\end{lstlisting}

验收标准：

\begin{itemize}[leftmargin=2em]
  \item Worker task 可以运行 golden example；
  \item \texttt{worker\_result.json} 生成；
  \item 输出缺失时 Worker 能报告失败；
  \item Worker 不做 fit selection 或 final physics decision；
  \item 测试通过。
\end{itemize}

\subsection{Week 4 任务 3：Month 1 final report}

学生需要写：

\begin{lstlisting}
docs/month1_final_report.md
\end{lstlisting}

报告内容：

\begin{enumerate}[leftmargin=2em]
  \item 项目目标理解；
  \item golden example 复现结果；
  \item data-flow map；
  \item config schema 设计；
  \item result schema 设计；
  \item data audit；
  \item pytest 测试列表；
  \item Worker wrapper 雏形；
  \item 当前限制；
  \item Month 2 计划。
\end{enumerate}

\subsection{Week 4 任务 4：Month 1 demo}

学生应准备一个 10--15 分钟展示，内容包括：

\begin{itemize}[leftmargin=2em]
  \item 一页说明项目目标；
  \item 一页说明 golden example 流程；
  \item 一页说明 repo 结构；
  \item 一页展示 CLI 运行命令；
  \item 一页展示 result.json / summary.csv / plots；
  \item 一页展示 pytest 结果；
  \item 一页说明 Worker wrapper；
  \item 一页说明当前问题和 Month 2 计划。
\end{itemize}

\section{第一个月导师验收清单}

导师可以按以下清单检查。

\subsection{复现性}

\begin{itemize}[leftmargin=2em]
  \item 是否能从 config 一键运行 golden example？
  \item 是否记录 git commit 和环境？
  \item 是否有 manifest.json？
  \item 是否有 run.log？
  \item README 是否足够另一个学生复现？
\end{itemize}

\subsection{代码质量}

\begin{itemize}[leftmargin=2em]
  \item 是否有清楚的目录结构？
  \item 是否把硬编码路径移入 config？
  \item 是否函数边界清楚？
  \item 是否避免大范围无关修改？
  \item 是否避免修改 raw data？
\end{itemize}

\subsection{测试}

\begin{itemize}[leftmargin=2em]
  \item pytest 是否通过？
  \item 测试是否覆盖 schema、data audit、CLI smoke test？
  \item 测试是否依赖机器特定路径？如果依赖，需要改掉。
\end{itemize}

\subsection{学生理解}

学生应能回答：

\begin{enumerate}[leftmargin=2em]
  \item 什么是 golden example？
  \item 为什么需要 config-driven pipeline？
  \item manifest.json 和 result.json 有什么区别？
  \item Data Auditor 为什么不能修改 raw data？
  \item Worker 为什么不能决定 central physics result？
  \item pytest 在科研分析中保护了什么？
  \item Codex 写的代码为什么不能直接相信？
\end{enumerate}

\section{第一个月常见错误}

\subsection{错误 1：一开始就让 Codex 写完整系统}

这会导致代码过大、不可控、学生无法理解。正确做法是每天一个小任务。

\subsection{错误 2：只跑 notebook，不记录命令}

科研 agent 系统必须可复现。没有命令记录和 config，就无法追责。

\subsection{错误 3：没有 schema}

没有 schema，config 和 result 会逐渐变成不稳定的自由字典，后续 agent 很难可靠读取。

\subsection{错误 4：没有 tests}

没有 tests，Codex 每次修改都可能破坏已有功能。

\subsection{错误 5：把 Worker 当成物理判断者}

Worker 可以运行和整理结果，但不能决定 central fit、丢弃数据或给出最终物理结论。

\subsection{错误 6：输出没有 provenance}

图和数值如果没有 config、git commit、run id 和日志，就很难用于严肃科研。

\section{Month 2 的衔接}

第一个月结束后，Month 2 将进入更高级阶段：

\begin{enumerate}[leftmargin=2em]
  \item 将 Month 1 的 CLI 包装进 Snakemake workflow；
  \item 增加 Fit Analyst 和 Skeptic Reviewer；
  \item 建立 fit candidate summary table；
  \item 实现 rule-based Supervisor；
  \item 写 \texttt{decision\_log.json}；
  \item 建立 accepted/rejected case memory。
\end{enumerate}

所以 Month 1 的关键不是把所有高级功能做完，而是打好以下基础：

\begin{lstlisting}
config -> schema -> CLI -> audit -> output -> tests -> Worker wrapper
\end{lstlisting}

只有这条链路清楚，Month 2 的 Snakemake 和 Supervisor 才能自然接上。

\section{附录：第一个月推荐 Codex Prompt 模板}

\subsection{通用模板}

\begin{lstlisting}
We are building an AI-assisted LQCD meson DA analysis-agent system.
Current stage: Month 1 foundation.

Task:
[describe one small task]

Files to edit:
[list files]

Requirements:
- [requirement 1]
- [requirement 2]
- [requirement 3]

Inputs:
[describe inputs]

Outputs:
[describe outputs]

Tests:
Please add pytest tests for:
- [test 1]
- [test 2]

Constraints:
- Do not modify raw data.
- Do not modify unrelated files.
- Do not make final physics decisions.
- Do not overwrite existing outputs unless explicitly allowed.
- Keep the implementation simple and readable.

After implementation, explain:
- what changed;
- how to run it;
- what tests were added.
\end{lstlisting}

\subsection{修 bug 模板}

\begin{lstlisting}
The following command failed:
[command]

Traceback:
[paste traceback]

Expected behavior:
[expected]

Actual behavior:
[actual]

Please fix only the cause of this failure.
Do not refactor unrelated code.
If appropriate, add a regression test so this failure is caught in the future.
\end{lstlisting}

\subsection{让 Codex 加测试的模板}

\begin{lstlisting}
Please add pytest tests for the existing implementation.

Focus on:
- valid input case;
- invalid input case;
- missing file or missing directory case;
- output schema validation;
- CLI smoke test using tmp_path.

Do not change the implementation unless a test reveals a real bug.
If you find a bug, explain it first and then make a minimal fix.
\end{lstlisting}

\section{第一个月最终标准}

第一个月完成后，学生应该达到以下水平：

\begin{itemize}[leftmargin=2em]
  \item 能复现导师给定的 golden example；
  \item 能将手动分析流程拆成输入、处理、输出和检查；
  \item 能写清楚 config schema 和 result schema；
  \item 能用 Codex 辅助生成代码，但知道如何检查和修正；
  \item 能写最小 pytest；
  \item 能用 Git 做小步提交；
  \item 能解释 Worker 的作用和边界；
  \item 能用一个命令运行基础 pipeline；
  \item 能提交一份清楚的 Month 1 final report。
\end{itemize}

如果学生完成以上内容，就已经具备继续进入 Month 2 高级 workflow 和 Supervisor agent 阶段的基础。

\end{document}
